% related.tex - Related Work

\section{Related Work}
\label{sec:related}

The most relevant prior work falls into several groups. %

\subsection{Capability Disciplines Without Static Tracking}

Capability systems from Dennis and Van Horn onward \cite{dennis-vanhorn}
established the core authority model: capabilities are explicit,
unforgeable references to authority. %
Miller's object-capability account
\cite{miller-robust,miller-paradigm} and languages such as E \cite{e-lang},
Joe-E \cite{joe-e}, and Newspeak \cite{newspeak} demonstrate that this
discipline scales to realistic programming languages. %

What these systems give is a strong semantic and engineering story about
authority flow. %
What they generally do not give is a small static account of
which capabilities a term may use, nor a compact proof-oriented model of
fine-grained call-site capability provision. %

\subsection{Capability-Based Module Systems}

Wyvern is a closely related language-based security design. %
Melicher \etal{} \cite{melicher-wyvern} propose a capability-based module
system for authority control, and Fish \etal{} \cite{fish-wyvern-io} present
a case study building an I/O library for Wyvern. %
\lambdaCap{} shares the goal that I/O authority should be represented by
capabilities rather than ambient effects, but studies a smaller formal core in
which capability access is reflected directly in term typing and call-site
capability provision. %

\subsection{Static Requirement Tracking}

Lucassen and Gifford \cite{lucassen-gifford} and Talpin and Jouvelot
\cite{talpin-jouvelot} introduced influential frameworks for tracking
program requirements in types. %
 \lambdaCap{} follows that general direction,
but applies it specifically to capability access. %

The key difference is not merely that \lambdaCap{} tracks capabilities. %
It is
that it tracks them in a setting where some capabilities remain under
call-site control rather than being uniformly captured or manually threaded. %
In that sense, the model is also close in spirit to coeffects: it describes
what a computation requires from its context \cite{petricek-coeffects}. %

\subsection{Verified Capability Reasoning}

Swasey \etal{} \cite{swasey-caps} verify object-capability patterns in a
richer language using logical relations. %
Their work is closer in proof style
than in modeling goal: they verify specific patterns in a realistic setting,
whereas \lambdaCap{} isolates a minimal core calculus with global capability
soundness. %

Devriese and Piessens \cite{devriese-effect} show, in a different security
setting, that static reasoning can establish strong semantic guarantees. %
That
result is not about capabilities, but it supports the general thesis that
small static models can carry real security meaning. %
